\section{Conclusion and Perspective}

\subsection{Summary of Achievements}

The project successfully implemented an RTOS-based smart climate control system that monitors temperature and humidity
and coordinates multiple actuator tasks using FreeRTOS-style asynchronous scheduling. The system fulfills the main
course requirements, including periodic sensor acquisition, LCD display updates, heater state indication, timed cooler
activation, and the multi-stage humidifier sequence. In addition, a dedicated Decision Maker task was introduced to
coordinate actuator execution and improve synchronization between monitoring and control tasks. This design reduces the
possibility of overlapping actuator activations and demonstrates the practical use of task coordination mechanisms in
real-time embedded systems.

\subsection{Limitations}

Several simplifications were made in this project. First, the heater, cooler, and humidifier are represented only by
RGB LED indicators rather than real physical actuators; therefore, the system does not actually modify the surrounding
temperature or humidity. Second, the project was developed and validated primarily in a simulation environment, so
additional testing on the YoluUNO hardware platform is necessary to evaluate timing accuracy, sensor reliability, and
overall system behavior under real operating conditions. Furthermore, the current control logic is threshold-based and
does not account for environmental dynamics such as delayed actuator effects or sensor noise.

\subsection{Future Work}

\begin{itemize}
  \item Implement the system on physical hardware using real actuators such as fans, heaters, and humidifiers.
  \item Port the software from Python to C++ to achieve better integration with embedded RTOS environments and more efficient
        low-level hardware control.
  \item Measure the response time of the physical actuators and incorporate this information into the control logic to prevent
        unnecessary repeated activations.
  \item Extend the system with closed-loop control methods such as PID control for smoother and more accurate regulation of
        temperature and humidity.
  \item Add features such as wireless monitoring, data logging, fault detection, and power-saving operating modes.
\end{itemize}